\clearpage\SourceLeaf{193}
\begin{center}
{\Large\textsc{Nachlass und Briefwechsel.}}
\end{center}

\bigskip
\begin{center}[ZUR PARALLELENTHEORIE.]\end{center}
\begin{center}\rule{2.5cm}{0.4pt}\end{center}

\begin{center}
\textsc{Gauss} an \textsc{Taurinus}. Göttingen, 8. November 1824.
\end{center}

\begin{center}Euer Wohlgeboren\end{center}

Ihr gefälliges Schreiben vom 30. Okt. nebst dem beigefügten kleinen Aufsatz habe ich nicht ohne
Vergnügen gelesen, umso mehr, da ich es sonst gewohnt bin, bei der Mehrzahl der Personen, die
neue Versuche über die sogenannte Theorie der Parallellinien [machen,] keine Spur von
wahrem geometrischem Geist anzutreffen. Gegen Ihren Versuch habe ich nichts (oder nicht viel)
anderes einzuwenden, als dass er unvollständig ist. Zwar lässt Ihre Darstellung des Beweises,
dass die Summe der drei Winkel eines ebenen Dreiecks nicht größer als $180^{0}$ sein kann,
in geometrischer Strenge noch zu wünschen übrig. Das ließe sich aber ergänzen,
und es besteht kein Zweifel, dass sich diese Unmöglichkeit mit äußerster
Strenge beweisen lässt. Ganz anders verhält es sich aber mit dem $2^{\mathrm n}.$ Teil,
dass die Winkelsumme nicht kleiner als $180^{0}$ sein kann. Das ist der eigentliche
Knoten, die Klippe, an der alles scheitert. Ich vermute, dass Sie sich noch nicht lange
